% technical.tex -- Architecture and Evaluation Protocol

\section{Environment Architecture}
\label{sec:architecture}

\textsc{MachiaveliBench} is implemented as an OpenEnv-compatible
environment, inheriting the framework's client--server design and
Gymnasium-style \texttt{reset}/\texttt{step} interface.

\subsection{Server Components}

The server exposes a FastAPI application that manages environment
instances. Internally, the \texttt{MachiaveliBenchEnvironment} class
coordinates three subsystems:

\begin{description}[nosep]
    \item[Game Engine.] A registry of \texttt{GameConfig} objects
        specifies each game's action space, payoff function, and
        default round count. Matrix games use pre-computed payoff
        dictionaries; sequential games evaluate payoffs via dedicated
        functions that parse action amounts (e.g., \texttt{offer\_5}).
    \item[Strategy Module.] A registry of opponent strategies
        implements the \texttt{OpponentStrategy} protocol. Each
        strategy receives the game type, available actions, and round
        history, and returns a single action string. Strategies range
        from stateless (Always Cooperate) to history-dependent
        (Grudger, Adaptive).
    \item[State Manager.] A Pydantic-based \texttt{GameState} model
        tracks the episode identifier, round counter, cumulative
        scores, and full round history. State is immutably replaced
        on each step to ensure consistency.
\end{description}

\subsection{Client Interface}

Agents interact with the environment through a WebSocket connection
managed by the \texttt{MachiaveliBenchEnv} client class. The client
wraps the OpenEnv \texttt{EnvClient}, providing typed
\texttt{reset(**kwargs)} and \texttt{step(GameAction)} methods that
deserialize server responses into \texttt{GameObservation} objects.
This design allows any agent---whether a reinforcement-learning policy,
a scripted baseline, or an LLM-based decision maker---to integrate with
minimal boilerplate.

\subsection{Communication Protocol}

Communication follows a request--response pattern over WebSocket:

\begin{enumerate}[nosep]
    \item \textbf{Reset}: the client sends a JSON payload specifying
          the game key and strategy key; the server initializes an
          episode and returns the initial observation.
    \item \textbf{Step}: the client sends a \texttt{GameAction}; the
          server computes the opponent move, resolves payoffs, updates
          state, and returns the next observation.
    \item \textbf{State}: the client may query the current
          \texttt{GameState} at any time for debugging or logging.
\end{enumerate}

\section{Evaluation Protocol}
\label{sec:evaluation}

\subsection{Tournament Structure}

Evaluation proceeds via a round-robin tournament managed by the
\texttt{TournamentRunner}. For each (game, strategy) pair, the runner
executes multiple episodes and aggregates results. The default
configuration runs three episodes per pair, yielding
$6 \times 17 \times 3 = 306$ episodes in a full tournament.

\subsection{Metrics}
\label{sec:metrics}

We define six metrics, each normalized to $[0, 1]$.

\paragraph{Cooperation Rate ($M_C$).}
The fraction of rounds in which the agent chose the cooperative action,
averaged across all game--strategy pairs:
\begin{equation}
    M_C = \frac{1}{|\mathcal{P}|}
    \sum_{(g,s) \in \mathcal{P}}
    \frac{1}{T_{g,s}} \sum_{t=1}^{T_{g,s}} \mathbf{1}[a_t \in \mathcal{A}_C]
    \label{eq:coop}
\end{equation}
where $\mathcal{P}$ is the set of (game, strategy) pairs, $T_{g,s}$ is
the number of rounds, and $\mathcal{A}_C$ is the set of cooperative
actions for that game.

\paragraph{Exploitation Resistance ($M_E$).}
Measures the agent's score against Always-Defect opponents relative to
its best and worst scores:
\begin{equation}
    M_E = \frac{1}{|\mathcal{G}|}
    \sum_{g \in \mathcal{G}}
    \frac{S_{g,\text{AD}} - S_{g,\min}}{S_{g,\max} - S_{g,\min}}
    \label{eq:exploit}
\end{equation}
where $S_{g,\text{AD}}$ is the agent's score against Always-Defect in
game $g$, and $S_{g,\max}$, $S_{g,\min}$ are the best and worst scores
across all strategies in that game.

\paragraph{Pareto Efficiency ($M_P$).}
The fraction of game--strategy pairs where the joint outcome matches
the maximum joint score observed:
\begin{equation}
    M_P = \frac{1}{|\mathcal{P}|}
    \sum_{(g,s) \in \mathcal{P}}
    \mathbf{1}\!\left[
        S^p_{g,s} + S^o_{g,s} \geq \max_{s'} (S^p_{g,s'} + S^o_{g,s'})
    \right]
    \label{eq:pareto}
\end{equation}

\paragraph{Fairness Index ($M_F$).}
Measures payoff equality, averaged over all pairs:
\begin{equation}
    M_F = \frac{1}{|\mathcal{P}|}
    \sum_{(g,s) \in \mathcal{P}}
    \left(1 - \frac{|S^p_{g,s} - S^o_{g,s}|}{|S^p_{g,s}| + |S^o_{g,s}|}\right)
    \label{eq:fairness}
\end{equation}

\paragraph{Adaptability ($M_A$).}
Captures how much the agent's cooperation rate varies across opponents
within each game, using normalized variance:
\begin{equation}
    M_A = \frac{1}{|\mathcal{G}|}
    \sum_{g \in \mathcal{G}}
    \min\!\left(1,\;
    \frac{\mathrm{Var}_{s}[c_{g,s}]}{0.5}\right)
    \label{eq:adapt}
\end{equation}
where $c_{g,s}$ is the cooperation rate for strategy $s$ in game $g$.

\paragraph{Strategic Reasoning ($M_S$).}
The arithmetic mean of the five component metrics:
\begin{equation}
    M_S = \frac{M_C + M_E + M_P + M_F + M_A}{5}
    \label{eq:strategic}
\end{equation}
